% !TEX program = xelatex
\documentclass[11pt,a4paper,fontset=fandol]{ctexart}

\newcommand{\coursecode}{SC2008}
\newcommand{\coursetitle}{Computer Networks}
\newcommand{\academicyear}{AY2026--27 Semester 1}
\usepackage{needspace}
\usepackage[margin=25mm]{geometry}
\usepackage{amsmath,amssymb,amsthm,bm,mathtools}
\usepackage{booktabs,tabularx}
\usepackage{enumitem}
\usepackage{graphicx}
\usepackage{xcolor}
\usepackage{tikz}
\usepackage{fancyhdr}
\usepackage{hyperref}
\usepackage[nameinlink,noabbrev]{cleveref}
\usetikzlibrary{arrows.meta,calc,fit,positioning}

\definecolor{ntuRed}{HTML}{D71440}
\definecolor{noteBlue}{HTML}{1F4E79}
\definecolor{verifyOrange}{HTML}{A64B00}

\hypersetup{
  colorlinks=true,
  linkcolor=noteBlue,
  citecolor=noteBlue,
  urlcolor=noteBlue,
  pdftitle={\coursecode: \coursetitle},
  pdfsubject={\academicyear\ 课程学习笔记}
}

\setlength{\parindent}{0pt}
\setlength{\parskip}{0.6em}
\setlist{itemsep=0.25em,topsep=0.35em}
\graphicspath{{figures/}}

\pagestyle{fancy}
\fancyhf{}
\lhead{\coursecode}
\rhead{\coursetitle}
\cfoot{\thepage}
\setlength{\headheight}{14pt}

\newtheorem{definition}{Definition（定义）}[section]
\newtheorem{theorem}[definition]{Theorem（定理）}
\newtheorem{proposition}[definition]{Proposition（命题）}
\newtheorem{example}[definition]{Example（例）}
\newtheorem{remark}[definition]{Remark（说明）}

\newcommand{\recordid}[1]{%
  \textcolor{noteBlue}{\textbf{编号：} \texttt{#1}}%
}

\newcommand{\lecturemeta}[6]{%
  \begin{center}
    \begin{tabularx}{\textwidth}{@{}lX@{}}
      \toprule
      课堂编号 & \texttt{#1} \\
      上课日期 & #2 \\
      主要资料 & #3 \\
      覆盖范围 & #4 \\
      视频补充 & #5 \\
      人工确认 & #6 \\
      \bottomrule
    \end{tabularx}
  \end{center}
}

\newcommand{\exercisemeta}[5]{%
  \begin{center}
    \begin{tabularx}{\textwidth}{@{}lX@{}}
      \toprule
      题组编号 & \texttt{#1} \\
      对应日期 & #2 \\
      主要来源 & #3 \\
      核对状态 & #4 \\
      人工确认 & #5 \\
      \bottomrule
    \end{tabularx}
  \end{center}
}

% Course introductions: lectures and tutorials use the same horizontal frame.
% Use noteintro when the existing source information is written as prose.
\newcommand{\noteintro}[1]{%
  \begin{center}
    \begin{tabularx}{\textwidth}{@{}X@{}}
      \toprule
      #1 \\
      \bottomrule
    \end{tabularx}
  \end{center}
}

% Keep the tutorial and lecture-example interfaces distinct. Written-assignment
% documents retain exercisemeta and their existing local formatting overrides.
\newcommand{\tutorialmeta}[5]{\exercisemeta{#1}{#2}{#3}{#4}{#5}}
\newcommand{\examplemeta}[5]{%
  \par\textbf{题组编号：}\texttt{#1}\quad\textbf{对应日期：}#2\par
  \textbf{主要来源：}#3\par
  \textbf{核对状态：}#4\quad\textbf{人工确认：}#5\par
}

\newcommand{\lecturetopic}[2]{%
  \subsection{#2}\label{topic:#1}%
  \recordid{#1}\par
}

% Course notes use category -> exercise group -> question. Written-assignment
% documents retain their existing section -> question hierarchy.
\newcommand{\exercisegroup}[2]{%
  \subsection{#2}\label{exercisegroup:#1}%
}

\newcommand{\exercisequestionheading}{\subsection}

\newenvironment{groupedexercises}
  {\renewcommand{\exercisequestionheading}{\subsubsection}}
  {}

\newcommand{\exercisequestion}[2]{%
  \exercisequestionheading{#2}\label{exercise:#1}%
  \recordid{#1}\par
}

% A compact notes TOC with enough number width for headings such as 14.10.3.
% The local group leaves written-assignment TOCs and body spacing unchanged.
\makeatletter
\newcommand{\coursenotestoc}{%
  \begingroup
  \small
  \setlength{\parskip}{0.15em}%
  \renewcommand*\l@subsection{\@dottedtocline{2}{1.5em}{3.0em}}%
  \renewcommand*\l@subsubsection{\@dottedtocline{3}{4.5em}{4.2em}}%
  \tableofcontents
  \endgroup
}
\makeatother

\newcommand{\relatedtopic}[2]{%
  \hyperref[topic:#1]{#2}~(\texttt{#1})%
}

\newcommand{\relatedexercise}[2]{%
  \hyperref[exercise:#1]{#2}~(\texttt{#1})%
}

\newcommand{\codeassignmentref}[2]{%
  \texttt{#1}~(\nolinkurl{#2})%
}

\newcommand{\verify}[1]{%
  \textcolor{verifyOrange}{\textbf{[Verify:} #1\textbf{]}}%
}

\newcommand{\keyidea}[1]{%
  \par\textcolor{noteBlue}{\textbf{核心。}} #1\par
}

\newenvironment{workedsolution}
  {\par\textbf{参考解答。}\par}
  {\hfill$\square$\par}

\newenvironment{checkedsolution}
  {\par\textbf{核对后的解答。}\par}
  {\hfill$\square$\par}

\title{\textcolor{ntuRed}{\coursecode}\\[0.4em]\coursetitle}
\author{课程学习笔记}
\date{\academicyear}

\linespread{1.12}
\setlength{\parskip}{0.25em plus 0.1em minus 0.05em}
\setlist{itemsep=0.1em,parsep=0pt,topsep=0.2em,partopsep=0pt}
\ctexset{
  section={beforeskip=1.2ex plus 0.3ex,afterskip=0.6ex},
  subsection={beforeskip=1.3ex plus 0.3ex,afterskip=0.5ex},
  subsubsection={beforeskip=1.1ex plus 0.2ex,afterskip=0.4ex}
}
\AddToHook{cmd/subsection/before}{\Needspace{6\baselineskip}}
\clubpenalty=10000
\widowpenalty=10000
\raggedbottom
\renewcommand{\exercisequestion}[2]{%
  \Needspace{5\baselineskip}%
  \subsubsection{#2}\label{exercise:#1}%
}
\renewcommand{\exercisemeta}[5]{%
  {\small\textbf{日期：}#2\quad\textbf{来源：}#3\par}%
}
\renewenvironment{checkedsolution}
  {\par\nopagebreak[4]\textbf{解答：}\ignorespaces}
  {\par}
\title{\textcolor{ntuRed}{\coursecode}\\[0.4em]\coursetitle\\[0.8em]\large 书面作业}
\author{作业整理}

\begin{document}
\begin{center}
  {\Large\bfseries\textcolor{ntuRed}{\coursecode}\quad\coursetitle\quad 书面作业\par}
  {\small\academicyear\par}
\end{center}

% Written assignment: SC2008-HW01
% Source: completed IP LAB 01 web activities 01--10, supplied in the conversation.
% Completion/check date: 2026-10-17; answers checked against page feedback.
\section{Online Interactive Exercise 1：IP Addresses}
\label{homework:SC2008-HW01}
\exercisemeta
  {SC2008-HW01}
  {2026-10-17}
  {IP LAB 01: Build a network that fits；网页 01--10}
  {练习与 checkpoint 已完成；依据题面、反馈及计算核对}
  {已确认提交}

本题组含七项 practice activities（练习活动）、两个网页例题及五道 checkpoint（综合检查）题。对应讲义 \texttt{SC2008-L06-T02}，与 Tutorial 05 Q1--Q3 的地址规划方法相同。

\subsection{计算规则}

IPv4 共 32 位，前缀 \texttt{/p} 留下 $h=32-p$ 个 host bits（主机位）。本练习均保留 network address（网络地址）与 broadcast address（广播地址），因此
\[
  B=2^{32-p}\quad\text{（总地址数）},\qquad U=B-2\quad\text{（可用接口数）}.
\]
已知所需接口数 $N$，取最小的 $h$ 使 $2^h-2\geq N$，再用前缀 $32-h$。题目只数电脑时需加路由器 LAN 接口；若总数已包含它，就不再增加。

Subnet mask（子网掩码）前 $p$ 位为 1、其余为 0。网络地址等于 IP 与掩码逐位 AND；广播地址保留前缀、将主机位全设为 1。可用范围去掉首末地址。
父块 \texttt{/p} 划为 \texttt{/q}（$q\geq p$）可得 $2^{q-p}$ 个等大小子网。VLSM（可变长子网掩码）允许各 LAN 使用不同前缀，需满足\textbf{容量足够、边界对齐、互不重叠、位于父块内}。

\subsection{练习活动}

\exercisequestion{SC2008-HW01-Q01}{Read the mask：前缀增长后的容量}

\textbf{题目。}主机地址固定为 \texttt{192.168.10.77}。将前缀从 \texttt{/26} 改为 \texttt{/27}，每块的\emph{总地址数}如何变化？选项：They stay the same（不变）、They halve（减半）、They double（翻倍）。

\begin{checkedsolution}
选 \textbf{They halve}。主机位从 6 位减到 5 位，总地址数从 $2^6=64$ 减到 $2^5=32$。
在 \texttt{192.168.10.0/24} 内，包含该主机的范围由 \texttt{.64--.127} 缩为 \texttt{.64--.95}。
\end{checkedsolution}

\exercisequestion{SC2008-HW01-Q02}{Find the range：求子网全部边界}

\textbf{题目。}给定 \texttt{172.20.6.174/27}，填写完整掩码、网络地址、广播地址、首末可用地址及可用地址数。网页先以 \texttt{172.20.6.141/27} 演示同一过程，二者对照如下。

\begin{checkedsolution}
\texttt{/27} 的块长为 32，最后一段起点为 $0,32,64,96,128,160,192,224$。
$\lfloor141/32\rfloor\times32=128$；$\lfloor174/32\rfloor\times32=160$。
\begin{center}
\small
\begin{tabular}{@{}lll@{}}
\toprule
字段 & 网页例题：\texttt{.141/27} & 本题：\texttt{.174/27} \\
\midrule
掩码 & \texttt{255.255.255.224} & \texttt{255.255.255.224} \\
网络地址 & \texttt{172.20.6.128} & \texttt{172.20.6.160} \\
广播地址 & \texttt{172.20.6.159} & \texttt{172.20.6.191} \\
首个可用地址 & \texttt{172.20.6.129} & \texttt{172.20.6.161} \\
最后可用地址 & \texttt{172.20.6.158} & \texttt{172.20.6.190} \\
可用地址数 & $32-2=30$ & $32-2=30$ \\
\bottomrule
\end{tabular}
\end{center}
掩码的最后一段为 $11100000_2=128+64+32=224$。已知主机 IP 用来定位所属块，不能把该 IP 当作子网起点往后数。
\end{checkedsolution}

\exercisequestion{SC2008-HW01-Q03}{Count every interface：是否已包含路由器}

\textbf{题目。}分别选择最小合适子网，选项均为 \texttt{/28、/27、/26、/25}。
Case A：30 台电脑，另需一个路由器 LAN 接口。Case B：总共需要 30 个可用地址，\emph{已包含}路由器 LAN 接口。

\begin{checkedsolution}
Case A 需要 $30+1=31$ 个可用地址，\texttt{/27} 只有 30 个，故选 \textbf{\texttt{/26}}（62 个可用）。
Case B 只需 30 个，选 \textbf{\texttt{/27}} 恰好满足。两题差别在计数口径。
\end{checkedsolution}

\exercisequestion{SC2008-HW01-Q04}{Split a parent block：等大小子网}

\textbf{题目。}父块为 \texttt{10.48.0.0/20}，六个部门各有至多 300 台电脑及一个网关接口。选最小的等大小子网（\texttt{/24、/23、/22、/21}），并求每部门可用地址数、父块可容纳的子网数、分配后剩余块数。

\begin{checkedsolution}
每部门需要 301 个地址；$2^8-2=254<301\leq2^9-2=510$，所以选 \textbf{\texttt{/23}}，每块 \textbf{510} 个可用地址。
父块能容纳 $2^{23-20}=\mathbf{8}$ 个子网，分配六块后剩 \textbf{2} 块。
\texttt{/23} 的掩码为 \texttt{255.255.254.0}，八块可从 \texttt{10.48.0.0、10.48.2.0、\ldots、10.48.14.0} 开始。
\end{checkedsolution}

\textbf{网页例题：不同大小的 LAN。}
在 \texttt{10.90.2.0/24} 中，X 有 60 台电脑及一个网关，Y 有 20 台电脑及一个网关。
X 需 61 个可用地址，取 \texttt{10.90.2.0/26}（\texttt{.0--.63}）；Y 需 21 个，取 \texttt{10.90.2.64/27}（\texttt{.64--.95}）。
Y 若取 \texttt{.32/27}，虽边界合法、容量足够，仍会与 X 重叠。

\exercisequestion{SC2008-HW01-Q05}{Build your address plan：四个 LAN 与四条链路}

\textbf{题目。}在 \texttt{10.60.4.0/24} 内规划 A、B、C、D 四个 LAN，电脑数依次为 90、45、24、10，\emph{每个 LAN 另需一个路由器接口}。还需四条 router-to-router links（路由器互连链路），每条用 \texttt{/30} 连接两个接口。为各 LAN 选前缀和合法网络地址，确保链路仍有空间。

\begin{checkedsolution}
先计接口，再按块大小从大到小放置。下表中的 \texttt{.x} 均表示 \texttt{10.60.4.x}。
\begin{center}
\small
\begin{tabular}{@{}lrrlll@{}}
\toprule
LAN & 需求 & 可用容量 & 网络地址/前缀 & 全范围 & 可用范围 \\
\midrule
A & 91 & 126 & \texttt{.0/25} & \texttt{.0--.127} & \texttt{.1--.126} \\
B & 46 & 62 & \texttt{.128/26} & \texttt{.128--.191} & \texttt{.129--.190} \\
C & 25 & 30 & \texttt{.192/27} & \texttt{.192--.223} & \texttt{.193--.222} \\
D & 11 & 14 & \texttt{.224/28} & \texttt{.224--.239} & \texttt{.225--.238} \\
\bottomrule
\end{tabular}
\end{center}
余下 \texttt{.240--.255} 共 16 个地址，恰好划成四个 \texttt{/30}。每条链路两端可分配：
\begin{center}
\small
\begin{tabular}{@{}llll@{}}
\toprule
链路 & 网络地址/前缀 & 两端接口地址 & 广播地址 \\
\midrule
1 & \texttt{.240/30} & \texttt{.241、.242} & \texttt{.243} \\
2 & \texttt{.244/30} & \texttt{.245、.246} & \texttt{.247} \\
3 & \texttt{.248/30} & \texttt{.249、.250} & \texttt{.251} \\
4 & \texttt{.252/30} & \texttt{.253、.254} & \texttt{.255} \\
\bottomrule
\end{tabular}
\end{center}
总占用 $128+64+32+16+4\times4=256$，正好用完父块。各块起点按自身块长对齐，且互不重叠；这是一套可行分配，并非唯一排列。
\end{checkedsolution}

\exercisequestion{SC2008-HW01-Q06}{Keep the given addresses：保留已有主机 IP}

\textbf{题目。}父块为 \texttt{10.70.6.0/24}。A 共需 30 个可用地址，含主机 \texttt{10.70.6.145}；D 共需 10 个，含主机 \texttt{10.70.6.174}。数量\emph{均已包含路由器接口}。从 \texttt{/26、/27、/28} 中选各自最小合适前缀，并给出网络地址。

\begin{checkedsolution}
A 取 \textbf{\texttt{10.70.6.128/27}}，可用范围 \texttt{.129--.158}，含 \texttt{.145}，提供 30 个可用地址。

D 取 \textbf{\texttt{10.70.6.160/28}}，可用范围 \texttt{.161--.174}，含 \texttt{.174}，提供 14 个。
两块不重叠。容量决定所需块大小，已有 IP 决定这个块应放在哪里；已有 IP 必须是可用地址，不能落在网络或广播地址上。
\end{checkedsolution}

\exercisequestion{SC2008-HW01-Q07}{Address the routers：逐接口分配地址}

\textbf{题目。}LAN A 为 \texttt{10.22.4.0/25}，已有主机使用 \texttt{10.22.4.1}。R1--R2 链路为 \texttt{10.22.4.240/30}。选择 R1 的 LAN 网关地址、R1 与 R2 的链路接口地址，以及两端链路掩码。

\begin{checkedsolution}
\begin{center}
\small
\begin{tabular}{@{}lll@{}}
\toprule
接口 & 一套可行地址 & 掩码 \\
\midrule
R1 LAN 网关 & \texttt{10.22.4.2} & \texttt{255.255.255.128} \\
R1 链路接口 & \texttt{10.22.4.241} & \texttt{255.255.255.252} \\
R2 链路接口 & \texttt{10.22.4.242} & \texttt{255.255.255.252} \\
\bottomrule
\end{tabular}
\end{center}
LAN 可用范围为 \texttt{.1--.126}，排除已占用的 \texttt{.1} 后可选 \texttt{.2}。链路只剩 \texttt{.241}、\texttt{.242} 两个可用地址，不能用网络地址 \texttt{.240} 或广播地址 \texttt{.243}。两个链路地址可以交换；同一路由器的不同接口各自使用所在子网的地址与掩码。
\end{checkedsolution}

\subsection{Independent checkpoint}

\exercisequestion{SC2008-HW01-Q08}{Question 1：Masks}
\textbf{题目。}What is the dotted-decimal mask for \texttt{/26}?
\begin{checkedsolution}
\textbf{\texttt{255.255.255.192}}。前 24 位为 1，最后一段为 $11000000_2=192$。
\end{checkedsolution}

\exercisequestion{SC2008-HW01-Q09}{Question 2：Address ranges}
\textbf{题目。}Host \texttt{192.168.40.201/27}: what is the subnet broadcast address?
\begin{checkedsolution}
块长 32，$\lfloor201/32\rfloor\times32=192$，范围为 \texttt{.192--.223}，故广播地址为 \textbf{\texttt{192.168.40.223}}。
\end{checkedsolution}

\exercisequestion{SC2008-HW01-Q10}{Question 3：Interface counting}
\textbf{题目。}A LAN needs 14 computers and one router interface. What is the smallest suitable prefix? 选项：\texttt{/25、/26、/27、/28、/29}。
\begin{checkedsolution}
选 \textbf{\texttt{/27}}。需 $14+1=15$ 个可用地址；\texttt{/28} 仅有 $16-2=14$ 个，\texttt{/27} 有 30 个。
\end{checkedsolution}

\exercisequestion{SC2008-HW01-Q11}{Question 4：Parent and child blocks}
\textbf{题目。}How many equal \texttt{/27} subnets fit inside one \texttt{/24} parent block?
\begin{checkedsolution}
\textbf{8 个}。前缀增加 3 位，故子网数为 $2^{27-24}=8$；也可由 $256/32=8$ 核对。
\end{checkedsolution}

\exercisequestion{SC2008-HW01-Q12}{Question 5：VLSM boundaries}
\textbf{题目。}LAN X uses \texttt{10.80.2.0/26}. LAN Y needs 25 usable addresses, including its router. Which proposed block for Y has enough capacity and no overlap?
\begin{enumerate}[label=\Alph*.,leftmargin=2.4em]
  \item \texttt{10.80.2.32/27}
  \item \texttt{10.80.2.64/28}
  \item \texttt{10.80.2.64/27}
  \item \texttt{10.80.2.0/25}
\end{enumerate}
\begin{checkedsolution}
选 \textbf{C：\texttt{10.80.2.64/27}}。X 占 \texttt{.0--.63}，C 占 \texttt{.64--.95}、提供 30 个可用地址。
A 占 \texttt{.32--.63}、D 占 \texttt{.0--.127}，均与 X 重叠；B 只有 14 个可用地址，不够 25 个。
\end{checkedsolution}

% Written assignment: SC2008-HW02
% Source: completed case #1042 web pages and simulated command outputs,
% supplied in the conversation; completion/check date: 2026-10-24.
\clearpage
\section{Online Interactive Exercise 2：Sales server unreachable}
\label{homework:SC2008-HW02}
\exercisemeta
  {SC2008-HW02}
  {2026-10-24}
  {交互练习案例 \#1042；题面、网络图及模拟命令输出}
  {已完成；依据页面反馈与地址计算核对}
  {已确认完成}

\subsection{故障场景与排查证据}

Engineering 员工 Mei Lin 的笔记本由承包商手动配置。她无法访问服务器 \texttt{192.168.40.70}，但打印和上网正常。以下为练习中的模拟网络与测试。

\texttt{ipconfig} 显示：笔记本 IP 为 \texttt{192.168.40.10}，掩码为 \texttt{255.255.255.0}（\texttt{/24}），默认网关为 \texttt{192.168.40.1}。

\begin{center}
\small
\begin{tabularx}{\textwidth}{@{}lX@{}}
\toprule
检查 & 页面显示的结果 \\
\midrule
\texttt{ping 192.168.40.70} & 四次 \texttt{Destination host unreachable}，消息来自笔记本自身 \texttt{192.168.40.10}。\\
\texttt{ping 192.168.40.20} & 打印机正常回复，四次均成功。\\
\texttt{ping 8.8.8.8} & 外网目标正常回复，四次均成功。\\
\texttt{arp -a} & 有网关 \texttt{.1} 和打印机 \texttt{.20} 的动态记录，没有服务器 \texttt{.70} 的记录。\\
\bottomrule
\end{tabularx}
\end{center}

第一条 ping 虽显示 \texttt{Received = 4, Lost = 0}，收到的仍是不可达消息，不能据此判断服务器连通。

\textbf{实际拓扑。}下图中 \texttt{.x} 均指 \texttt{192.168.40.x}；Router R 通过三个接口连接各 LAN。
\begin{center}
\begin{tikzpicture}[
  lan/.style={draw=noteBlue,rounded corners=2pt,text width=4.25cm,align=center,inner sep=5pt,font=\small},
  port/.style={font=\small,fill=white,inner sep=2pt},
  line/.style={thick,draw=noteBlue}
]
  \node[lan] (eng) at (0,0) {\textbf{Engineering LAN}\\\texttt{192.168.40.0/26}\\范围 \texttt{.0--.63}\\笔记本 \texttt{.10}（误设 \texttt{/24}）\\打印机 \texttt{.20}；DNS \texttt{.2}};
  \node[lan] (sales) at (5.15,0) {\textbf{Sales LAN}\\\texttt{192.168.40.64/27}\\范围 \texttt{.64--.95}\\服务器 \texttt{.70}\\Sales PC \texttt{.80}};
  \node[lan] (support) at (10.3,0) {\textbf{Support LAN}\\\texttt{192.168.40.96/28}\\范围 \texttt{.96--.111}\\Support PC \texttt{.100}\\\strut};
  \node[draw=noteBlue,rounded corners=2pt,inner sep=6pt,font=\small\bfseries] (router) at (5.15,-2.1) {Router R};
  \draw[line] (eng.south) |- (router.west);
  \draw[line] (sales.south) -- (router.north);
  \draw[line] (support.south) |- (router.east);
  \node[port] at (1.35,-2.1) {接口 \texttt{.1}};
  \node[port,right] at (5.15,-1.5) {接口 \texttt{.65}};
  \node[port] at (8.9,-2.1) {接口 \texttt{.97}};
\end{tikzpicture}
\end{center}

\subsection{诊断与修复}

\exercisequestion{SC2008-HW02-Q01}{先做测试，再判断故障位置}
\textbf{题目。}面对上述故障，如何回应？选项：Let me run a few tests on your laptop；The Sales server must be down。
\begin{checkedsolution}
选 \textbf{Let me run a few tests on your laptop}。客户端配置也可能导致故障，应先检查配置与连通性。
\end{checkedsolution}

\exercisequestion{SC2008-HW02-Q02}{当前掩码算出哪个网络地址？}
\textbf{题目。}服务器 \texttt{192.168.40.70} 按当前 \texttt{/24} 掩码计算，网络地址是哪个？选项：\texttt{192.168.40.70}、\texttt{192.168.40.64}、\texttt{192.168.40.0}。
\begin{checkedsolution}
选 \textbf{\texttt{192.168.40.0}}。逐位 AND 后，前三段保留，最后一段清零；笔记本 \texttt{.10} 也得到同一网络地址。
\end{checkedsolution}

\exercisequestion{SC2008-HW02-Q03}{笔记本怎样决定发送路径？}
\textbf{题目。}根据上述计算，笔记本会作何判断？选项：\texttt{.70} is remote, send it to the gateway；\texttt{.70} is local, send it directly。
\begin{checkedsolution}
选 \textbf{local，直接发送}。错误的 \texttt{/24} 让笔记本把 \texttt{.70} 判断为本地地址，尝试直接取得服务器 MAC。
\end{checkedsolution}

\exercisequestion{SC2008-HW02-Q04}{服务器实际上在哪个 LAN？}
\textbf{题目。}在 Engineering、Sales、Support 三个 LAN 中，选择 \texttt{192.168.40.70} 实际所在的 LAN。
\begin{checkedsolution}
选 \textbf{Sales LAN：\texttt{192.168.40.64/27}}。\texttt{.70} 属于 \texttt{.64--.95}，不在 Engineering 的 \texttt{.0--.63} 内，访问它需经过 Router R。
\end{checkedsolution}

\exercisequestion{SC2008-HW02-Q05}{修复前，ARP 查询哪个 IP？}
\textbf{题目。}笔记本误认为 \texttt{.70} 在本地，发送前会用 ARP 查询哪个 IP 的 MAC？选项如下：
\begin{enumerate}[label=\Alph*.,leftmargin=2.4em]
  \item \texttt{255.255.255.255}（所有设备）；
  \item \texttt{192.168.40.70}（服务器）；
  \item \texttt{192.168.40.65}（Router R 的 Sales 接口）；
  \item \texttt{192.168.40.1}（本机默认网关）。
\end{enumerate}
\begin{checkedsolution}
选 \textbf{B：\texttt{192.168.40.70}}。ARP 查询的目标 IP 为 \texttt{.70}，以太网目的 MAC 为广播地址 \texttt{ff:ff:ff:ff:ff:ff}。按本题普通路由转发设置，ARP 广播只在 Engineering LAN 内传播，无法到达 Sales 服务器；所以笔记本无法解析服务器 MAC，报告不可达。
\end{checkedsolution}

\exercisequestion{SC2008-HW02-Q06}{修正 IPv4 配置并验证}
\textbf{题目。}在 IPv4 Properties 中手动填写 IP、掩码、默认网关及首选 DNS，然后再次测试服务器。
\begin{checkedsolution}
根据笔记本所在的 \textbf{Engineering LAN} 填写：
\begin{center}
\small
\begin{tabular}{@{}ll@{}}
\toprule
字段 & 正确值 \\
\midrule
IP address & \texttt{192.168.40.10} \\
Subnet mask & \textbf{\texttt{255.255.255.192}}（\texttt{/26}）\\
Default gateway & \texttt{192.168.40.1} \\
Preferred DNS server & \texttt{192.168.40.2} \\
\bottomrule
\end{tabular}
\end{center}
故障修正点是掩码，IP 与网关保持不变。DNS 来自拓扑图；直接 ping 数字 IP 无需名称解析。

比较最后一段：$10\mathbin{\&}192=0$，$70\mathbin{\&}192=64$。结果不同，笔记本将发往 \texttt{.70} 的包交给网关 \texttt{.1}。这只判断目标是否属于本地子网；Sales 的实际前缀仍是图中的 \texttt{/27}。

重新运行 \texttt{ipconfig} 核对配置，再运行 \texttt{ping 192.168.40.70}；连通的判据是收到来自 \texttt{.70} 的 Echo Reply。窗口里的旧输出不会随配置修改而刷新。
\end{checkedsolution}

\exercisequestion{SC2008-HW02-Q07}{修复后，首跳数据帧的目的 MAC}
\textbf{题目。}修复后，笔记本承载 IP 包的数据帧，其目的 MAC 属于谁？选项：The Sales server；Router R (its \texttt{.1} interface)；Broadcast \texttt{ff-ff-ff-ff-ff-ff}。
\begin{checkedsolution}
选 \textbf{Router R 的 \texttt{.1} 接口}，已有 ARP 表给出其 MAC 为 \texttt{00-1b-54-a3-20-01}。题目问的是承载 IP 包的数据帧；若无网关 MAC 缓存，发送前另需 ARP 查询。
\end{checkedsolution}

\exercisequestion{SC2008-HW02-Q08}{帧内 IP 包的目的地址}
\textbf{题目。}帧内 IP 包的目的 IP 是哪个？选项：\texttt{192.168.40.70}（服务器）；\texttt{192.168.40.1}（网关）。
\begin{checkedsolution}
选 \textbf{\texttt{192.168.40.70}}。\textbf{目的 MAC 指向当前链路的下一跳，目的 IP 指向最终服务器。}Router R 换用 Sales 链路的 MAC 地址封装并转发；本题无 NAT，IP 目的地址保持不变。
\end{checkedsolution}

\end{document}
